\subsection{ICG de resultados}

La primera implicancia consiste en preservar el ICG como una señal interpretable del desempeño alcanzado por la empresa prestadora. Como se mostró en la Sección~3.1 y se formalizó en la Sección~4.2, la incorporación de variables de ejecución reduce la sensibilidad del índice frente a los resultados y dificulta interpretar su valor agregado. La magnitud práctica de este efecto se evalúa mediante la simulación contrafactual de la Sección~5.

Por ello, el ICG debería concentrarse en variables que representen condiciones efectivamente alcanzadas en la prestación, la operación o el desempeño económico-financiero, y no en los medios utilizados para producirlas. La ejecución de una inversión o la implementación de un instrumento de gestión continuaría siendo exigible, pero solo afectaría el ICG cuando se traduzca en una mejora observable del resultado correspondiente.

La selección de los indicadores debería considerar, al menos, cuatro criterios: que la variable represente una condición alcanzada y no únicamente la ejecución de una actividad; que pueda verificarse mediante información confiable y trazable; que presente un grado razonable de controlabilidad por parte de la empresa; y que sea relevante para la calidad, eficiencia o sostenibilidad de la prestación. Estos criterios permiten incluir tanto resultados finales, como continuidad, presión o acceso, como resultados operativos o económico-financieros, siempre que se mantenga una diferenciación clara respecto de las metas que miden actividades, instrumentos o uso de recursos.

La concentración del ICG en resultados no requiere eliminar la evaluación por localidad ni los umbrales individuales. Estos mecanismos pueden mantenerse para identificar brechas territoriales y evitar que un desempeño favorable en determinadas dimensiones compense niveles inaceptables en otras. Asimismo, el rediseño debería revisar la ponderación de los indicadores, de modo que responda a su importancia regulatoria y no únicamente al número de metas incorporadas en cada estudio tarifario.

